\documentclass{article}
\usepackage[T2A]{fontenc}
\usepackage[utf8]{inputenc}   
\usepackage[english, russian]{babel}

% Set page size and margins
% Replace `letterpaper' with`a4paper' for UK/EU standard size
\usepackage[a4paper,top=2cm,bottom=2cm,left=2cm,right=2cm,marginparwidth=1.75cm]{geometry}

\usepackage{amsmath}
\usepackage{graphicx}
\usepackage[colorlinks=true, allcolors=blue]{hyperref}
\usepackage{amsfonts}
\usepackage{amssymb}
% \usepackage[left=1cm,right=1cm,top=1cm,bottom=1cm]{geometry}
\usepackage{hyperref}
\usepackage{seqsplit}
\usepackage[dvipsnames]{xcolor}
\usepackage{enumitem}
\usepackage{algorithm}
\usepackage{algpseudocode}
\usepackage{algorithmicx}
\usepackage{mathalfa}
\usepackage{mathrsfs}
\usepackage{dsfont}
\usepackage{caption,subcaption}
\usepackage{wrapfig}
\usepackage[stable]{footmisc}
\usepackage{indentfirst}
\usepackage{rotating}
\usepackage{pdflscape}

\usepackage{minted}

\usepackage{MnSymbol,wasysym}

\title{Домашнее задание 1}
\author{Лиза Фоменко}
\date{}

\begin{document}

\maketitle



\section*{\textbf{Задание 0}}

\textbf{\textit{Задание с лекции: $\log_5 n =? O(\log_2 n)$}}

Вспомним, что  $\log_2 n = \frac{\log_5 n}{\log_5 2}$ и воспользуемся определеним O-нотации:
$\exists C > 0 : \exists N \in \mathbb{N} : \forall n \in \mathbb{N}, n \geq N$  верно, что $f(n) \leq C * g(n)$. 

Подставим уравнение в определение:
$$\log_5 n \leq C * \log_2 n = C * \frac{\log_5 n}{\log_5 2}$$

Видно, что при $C \geq \log_5 2$ равенство выполняется, откуда $\log_5 n = O(\log_2 n)$

(\textit{верно для любых логарифмов})

\bigskip

\section*{\textbf{Задание 1}}

\textbf{\textit{Известно, что $f(n) = O(n^2),\ g(n) = \Omega(1),\ g(n) = O(n)$. Положим $$h(n) = \frac{f(n)}{g(n)}.$$}}

\medskip

\textit{\textbf{1.} Возможно ли, что \textbf{а)} $h(n) = \Theta(n\log n)$; \textbf{б)} $h(n) = \Theta(n^3)$? Если да, приведите конкретные фунцкии. Если нет, докажите, что это невозможно.}

Так как у нас нет оценки $\Omega(f(n))$, то мы можем оценить только верхнюю границу функции $h(n)$. 

$$h(n) = \frac{f(n)}{g(n)} \leq \frac{O(f(n))}{\Omega(g(n))} = \frac{C * {n^2}}{C'} = C'' * n^2$$

Данным действием мы ограничили $O(h(n))$ сверху. $O(h(n))$ не может быть ассимптотически больше, чем $n^2$, но может быть меньше нее. Отсюда ответ на пукнт б.

\smallskip

\textbf{пункт б:} нет, такое не возможно. По определению $\Theta(h(n))$, $h(n)$ должна быть $\Omega(n^3)$, что невозможно, так как верхняя оценка сложности составляет $O(n^2)$ (более формально можно доказать так: $h(n) \leq C*n^2$; $n^2 < n^3$ для $\forall n > 1 \Rightarrow h(n) < C*n^3$, а по определению $\Omega$ должно быть $h(n) \geq C*n^3$)

\smallskip

\textbf{пункт а:} Здесь можно воспользоваться тривиальным свойством (так называемое "свойство рефлексивности" в Кормене): $h(n) = \Theta(h(n))$.
Тогда, зная, что $\Theta(h(n)) = n\log n$, подберем $f(n), g(n)$ (сохраняя оценки) так, чтобы $h(n) = \frac{f(n)}{g(n)} = n\log n$. 

При $f(n) = n^2$, $g(n) = \frac{n}{\log n}$, $h(n) = \frac{f(n)}{g(n)} = \frac{n^2}{\frac{n}{\log n}} = n\log n = \Theta(n\log n)$

\medskip

\textit{\textbf{2.} Приведите наилучшие (из возможных) верхние и нижние оценки на функцию $h(n)$ и приведите пример функций $f(n)$ и $g(n)$, для которых ваши оценки на $h(n)$ достигаются.}

\medskip

$\begin{cases}
  \frac{\Omega(f(n))}{n} \leq  \Omega(f(n)) \leq \frac{f(n)}{g(n)} \leq n \leq n^2 & \text{    if       } \Omega(f(n)) \leq n \\
  $$\frac{\Omega(f(n))}{n} \leq  n \leq \frac{f(n)}{g(n)} \leq \Omega(f(n)) \leq n^2$$  & \text{    if       } n \leq \Omega(f(n)) \leq n^2
\end{cases}$

\medskip


Нас просят найти наибольшую нижнюю оценку и наименьшую верхнюю границу для функции $h(n)$ в виде $n^a\log^bn$. Введем следующие предположения$^*$, которые не прописаны в задаче, но, кажется, предполагаются:

\begin{enumerate}
  \item $f(n) = O(n^2) \Rightarrow$ $n^2$ -- это минимальная степень полинома, для которого выполняется оценка $f(n) = O(n^x)$ (то есть $f(n) \neq O(n)$, например).
  
  \item Аналогично $g(n) = O(n) \Rightarrow$ $n$ -- это минимальная степень полинома, для которого выполняется оценка $f(n) = O(n^x)$ (то есть $f(n) \neq O(1)$, например).
\end{enumerate}

\medskip

Мы точно знаем, что $g(n) \leq C*n$, то есть имеет вид $\frac{n}{\log^an}, a \geq 0$ или $\log^bn, c \in \mathbb{R}$. 

\smallskip

В случае $f(n)$ функция имеет общий вид $n\log^an, a \geq 0$ или $\frac{n^2}{\log^bn}, b \geq 0$. Все перечисленные функции являются $O(n^2)$. 

\smallskip

Тогда из кобинаций всех значений $g(n)$ и $f(n)$ находим ту, которая имеет наименьшую $\Theta$ ассимптотику, и наибольшую $\Theta$ ассимптотику (почему? потому что все возможные значения функции $h(n)$ в такой постановке задачи выглядят как $n^a\log^bn$, а именно в этом классе функций мы ищем оценки. 

\smallskip

\begin{enumerate}
  \item $\frac{n^2}{\log^\alpha n}$ (при $f(n) = \frac{n^2}{\log^\alpha n}$, $g(n) = \log^bn$, $a \geq b$)
  
  \item $n\log^\alpha n$ (при $f(n) = \frac{n^2}{\log^\alpha n}$, $g(n) = \frac{n}{\log^b n}$ или $f(n) = n\log^\alpha n$, $g(n) = \log^bn$)
  
  \item  $\log^\alpha n$ (при $f(n) =n\log^an$, $g(n) = \frac{n}{\log^an}$)
\end{enumerate}

$$h(n) = O(n\log^\alpha n), \alpha \geq 0$$ 
$$h(n) = \Omega(\log^\alpha n), \alpha \geq 0$$

* А что, если не использовать предположения 1-2, то есть $f(n)$ может быть $O(n)$? Тогда для $f(n)$ будут справедливы те же оценки, что и для $g(n)$. Тогда при $g(n) = f(n)$ получаем $h(n) = 1 = \Omega(1)$. Верхняя оценка останется аналогичной, то есть $O(n\log^\alpha n), \alpha \geq 0$.

\section*{\textbf{Задание 2}}

\textbf{\textit{Найдите $\Theta$-асимптотику $\sum\limits_{i=1}^n \sqrt{i^3+2i+5}$.}}

Оценим каждое слагаемое суммы:

$$i^{3/2}= \sqrt{i^3} \leq \sqrt{i^3+2i+5} \leq \sqrt{i^3 + 2i^3 + 5i^3} = \sqrt{8i^3} = 2\sqrt{2}i^{3/2}$$

Тогда для оценки суммы верно, что 
$$\sum\limits_{i=1}^n i^{3/2} \leq \sum\limits_{i=1}^n \sqrt{i^3+2i+5} \leq 2\sqrt{2}\sum\limits_{i=1}^n i^{3/2}$$

Тогда, используя теорему из задания 3 ($\alpha = 3/2$)

$$\Omega(n^{5/2}) = \sum\limits_{i=1}^n i^{3/2} \leq \sum\limits_{i=1}^n \sqrt{i^3+2i+5} \leq 2\sqrt{2}\sum\limits_{i=1}^n i^{3/2} = O(n^{5/2})$$

откуда, по определению $\Theta$, 
$$\sum\limits_{i=1}^n \sqrt{i^3+2i+5} = \Theta(n^{5/2})$$
\bigskip

\section*{\textbf{Задание 3}}

\textbf{\textit{Докажите, не используя интегрального исчисления, что асимптотика $\sum^n_{i=1} i^{\alpha}
= \Theta(n^{1+\alpha})$, если
$\alpha > 0$.}}

$$\sum\limits_{i=1}^n i^\alpha = 1^\alpha + 2^\alpha + ... + n^\alpha \leq n^\alpha + n^\alpha + ... + ^\alpha = n*n^\alpha = n^{1+\alpha} = O(n^{1+\alpha})$$ 
\medskip
$$\sum\limits_{i=1}^n i^\alpha \geq \sum\limits_{i=n/2}^n i^\alpha = \frac{n^\alpha}{2^\alpha} + \frac{(n+2)^\alpha}{2^\alpha} + ... + n^\alpha \geq \frac{n}{2} * \frac{n^\alpha}{2^\alpha} = \frac{n^{1 + \alpha}}{2^{1 + \alpha}} = \Omega(n^{1+\alpha})$$ 

откуда, по определению $\Theta$, верно, что 
 $$\sum\limits_{i=1}^n i^\alpha = \Theta(n^{1+\alpha})$$

\textit{(доказательство для $\Omega$ можно проводить, в принципе, для любого $k \leq n$, я использовала $k = 2$)} 
\bigskip

\section*{\textbf{Задание 4}}

\textbf{\textit{Найдите $\Theta$-асимптотику функции
$g(n) = \sum ^n_{k=1} \frac{1}{k};$}}

$$\sum ^n_{k=1} \frac{1}{k} \geq \int_1^{n+1} \frac{dx}{x}= \ln(n + 1) \geq \ln(n) \Rightarrow  \Omega(\ln n)$$

$$\sum ^n_{k=1} \frac{1}{k} \leq 1 + \int_1^n \frac{dx}{x}= 1 + \ln(n) \Rightarrow  O(\ln n)$$

тогда, по определению $\Theta$, верно, что 
$$\sum ^n_{k=1} \frac{1}{k} = \Theta(\ln n)$$

\end{document}